\documentclass{article}
\usepackage[utf8]{inputenc}
\usepackage{enumitem}
\usepackage{dsfont}
\usepackage[margin=1in]{geometry}

\title{Placement progress and notes}
\author{Madeleine Watson}

\begin{document}

\maketitle

\tableofcontents

%---------
% MONTH 1
%---------

\section{Month 1: 23/02/2026 - 20/03/2026}
\subsection{Overview}

Spent the first week familiarising myself with beam optics concepts and Xsuite functions. Then began understanding and editing a code for a compact ring with 90$^\circ$ phase advance.

\subsection{Work completed}

\begin{itemize}
\item Initially read papers to understand how the ring functions and what processes go into it.
\item Made notes on initial ring design code and started compiling work into a GitHub repository.
\item Added sextupoles to the ring design and performed chromaticity correction.
\item Corrected beta-functions to create uniform Twiss plot.
\item Corrected working point and found two potential candidates.
\end{itemize}

\subsection{Problems encountered}

\textbf{Breaking the ring:} \texttt{Invalid n1}, \texttt{Invalid n2}, and \texttt{Invalid n3} errors arose when trying to add sextupoles. This scenario halted progress the most.
\begin{itemize}
    \item \textbf{Solution:} There was no set solution to this error but there were some common causes: changing the geometry of the ring, not matching/rematching the quadrupole strengths. The `breakage' could be identified by computing a Twiss with specified conditions.
\end{itemize}

\textbf{Sextupole positioning:} This would cause ring breakage. It took a week or so to identify sextupole positions that worked.
\begin{itemize}
    \item \textbf{Solution:} I tried creating a function that identified the best positions for sextupoles, but the best solution was trial and error and knowledge of ring properties.
\end{itemize}

\textbf{Large beta-functions:} Changing parameters resulted in increased beta-functions
\begin{itemize}
    \item \textbf{Solution:} Consistently running matching functions and identifying best positions and working points.
\end{itemize}

\subsection{Results and progress}
\begin{itemize}
    \item The initial compact ring now has sextupoles inserted: there are currently 16 in each arc (8 focusing, 8 defocusing).
    \item Strengths and beta-functions have been corrected
    \item Two potential working points have been identified and plotted on resonance graphs to show suitability
    \item Using these working points resulted in reduced beta-functions, more uniform Twiss plots, and slightly reduced sextupole strengths.
    \item Another configuration with a decreased triplet drift length is being explored and shows signs of reduced beta-functions
    \item Chromatic functions are plotted and show slightly different, but periodic, results for both working points.
\end{itemize}

\subsection{Next steps}
\begin{itemize}
    \item Further refine suitability of the working point
    \item Track off-momentum particles
    \item Track dynamic aperture and momentum acceptance
    \item Work on different sextupole configurations with different positions, families, and cell lengths in parallel
\end{itemize}

\subsection{Skills developed}
\begin{itemize}
    \item \textbf{Knowledge of beam optics concepts:} Twiss parameters, chromaticity, tune, basics of synchrotron radiation, and dynamic aperture
    \item \textbf{Xsuite:} Building rings and inserting elements, calculating Twiss plots, creating matching functions
\end{itemize}

%---------
% MONTH 2
%---------

\section{Month 2: 23/03/2026 - 24/04/2026}
\subsection{Overview}

Completed dynamic aperture and momentum acceptance studies for different configurations. Experimented with slicing, integrators, and fringe fields. Now introducing misalignments and corrections.

\subsection{Work completed}

\begin{itemize}
\item Plotted DA/MA for 7 different sextupole configurations 
\item Optimised configurations to increase DA/MA
\item Changed number of slices in integrators to increase quality of plots
\item Introduced fringe fields for every magnet type and compared with "perfect" results
\item Performed frequency map analysis
\item Introduced misalignments, inserted BPMs and correctors, and performed orbit correction
\end{itemize}

\subsection{Problems encountered}

\textbf{Reducing non-linearities:} Chromatic correction is working for $dqx$ and $dqy$ but often plotting tune against momentum deviation showed that $ddqx$ and $ddqy$ had not been properly corrected. This contributed to smaller DA/MA. Performing chromaticity correction to get these values to 0 often failed as it was too much of a jump.
\begin{itemize}
    \item \textbf{Solution:} The chromaticity correction tries to match to 80\% of the value of $ddqx$ and $ddqy$. If this is successful, it iterates through 80\% of the resulting values until the code fails. It takes the strengths from the iteration with the smallest second order chromaticity.
\end{itemize}


\textbf{Code organisation:} Many scripts and results in different places
\begin{itemize}
    \item \textbf{Solution:} Everything has been modularised so it is easier to see where errors arise and can edit things more efficiently. This will also help with later lattice designs.
\end{itemize}

\subsection{Results and progress}
\begin{itemize}
    \item Increasing the number of slices for the quadrupoles created smoother curves in the DA/MA plots, but there were no differences when increasing the slices for bends and quadrupoles.
    \item Fringe fields reduced the less stable configurations' DA/MA considerably, but not as much for stable configurations.
    \item Added another focusing sextupole pair to the configuration with two sextupole pairs at 180$^\circ$ phase advance, which dramatically improved momentum acceptance.
    \item Frequency map analysis calculated and shows stable diffusion
    \item BPMs and correctors are currently placed around every quadrupole (horizontal and vertical) and the MICADO algorithm is used for corrections
    \item With corrections, the orbits are currently oscillating at less than 1mm
\end{itemize}

\subsection{Next steps}
\begin{itemize}
    \item Studies on spin tracking and polarization build-up time
    \item Looking at integration with the rest of the injector complex using current macroparticle distribution
    \item Assess injection efficiency
\end{itemize}

\subsection{Skills developed}
\begin{itemize}
    \item \textbf{Dynamic aperture and momentum acceptance:} how it is calculated and the effects of adjusting integrators and slices
    \item \textbf{Knowledge of beam optics concepts:} fringe fields, FMA, BPMs, correctors, synchrotron radiation damping
\end{itemize}

%---------
% MONTH 3
%---------

\section{Month 3: 27/04/2026 - 22/05/2026}
\subsection{Overview}
Simulated injection efficiency using the most recent macroparticle distribution and varied reference energy. Observed the effect of misalignments on dynamic aperture. Created a code based on A.Verdier's paper that identifies dangerous resonances. Spent time studying theory in a more in depth way.
\subsection{Work completed}

\begin{itemize}
\item Using the most recent macroparticle distribution:
    \begin{itemize}
    \item Plotted transverse and longitudinal phase space
    \item Calculated emittances, Twiss parameters, and momentum spread
    \item Normalised distribution (taking dispersion into account) and performed initial injection simulation 
    \item Performed energy scan to find most efficient injection energy and ran simulation at this value 
    \end{itemize}
\item Ran dynamic aperture calculation with misalignments and corrections to see how much it is affected (for 2 configurations)
\item Created code that identifies which resonances the working point could potentially be excited to, for a specific lattice
\item Spent time studying the derivation of particle motion equations, limits of emittance, and limits of synchrotron radiation
\end{itemize}

\subsection{Problems encountered}

\textbf{Macroparticle distribution size:} The distribution contained 25,000 macroparticles which took up too much memory to be able to run in one go.
\begin{itemize}
    \item \textbf{Solution:} Took random subsets of 500 particles and averaged across seeds to get the injection efficiency.
\end{itemize}

\subsection{Results and progress}
\begin{itemize}
    \item Using configuration 1 at the average energy obtained $\approx 33\%$ injection efficiency
    \item Using configuration 2 at the average energy obtained $\approx 57\%$ injection efficiency
    \item Optimising the energy for injection increased the injection efficiency for configuration 2
    \item For both configurations, DA was around $5\sigma$ after misalignments and corrections but MA significantly reduced for configuration 2
    \item The Verdier resonance scan revealed that the current working point is close to several non-structural resonances
\end{itemize}

\subsection{Next steps}
\begin{itemize}
    \item Run DA/MA for newly calculated emittances 
    \item Simulate an energy compressor to see if the injection efficiency can increase further and reduce energy spread
    \item Spin tracking and equilibrium polarization
\end{itemize}

\subsection{Skills developed}
\begin{itemize}
    \item Simulations of particle distributions in the ring 
    \item Particle motion understanding
\end{itemize}

%---------
% MONTH 4
%---------

\section{Month 4: 26/05/2026 - 03/07/2026}
\subsection{Overview}
Used an energy compressor with optimised parameters to reduce energy spread and increase injection efficiency. Coded spin tracking simulations to calculate equilibrium polarization for multiple misalignment and correction cases.

\subsection{Work completed}

\begin{itemize}
\item Re-worked code structure to make it more sreamlined and ran configurations with newly calculated emittances
\item Introduced an energy compression system into the macroparticle script distribution using $\delta_{new}=\delta_{old}+V_{deb}\sin{2\pi\cdot\zeta_{new}+\frac{V_{RF}}{c}+\phi_{deb}}$ where $\zeta_{new}=\zeta_{old}+R_{56}\cdot\delta$
\item $V_{deb}$,$R_{56}$, and $\phi_{deb}$ were optimised to minimise energy spread
\item Studied theory for Thomas-BMT equation, spin tune and motion, Debrenev-Kondratenko-Mane and Baier-Khatkov-Strakhovenko formalisms
\item Built a spin tracking code based on this theory that calculates $P_{eq}$ for multiple misalignment seeds
\end{itemize}

\subsection{Problems encountered}

\textbf{Very low equilibrium polarization values:} Each result produced very low equilibrium polarization values which only became worse when corrections were implemented. However, there was no obvious error with the calculations.
\begin{itemize}
    \item \textbf{Solution:} The anomolous magnetic moment was not explicitly set which greatly affected results. Once this was corrected, results dramatically improved. To confirm the magnetic moment was correct; a plot of the invariant spin vector, using a particle with initialised spin, was used.
\end{itemize}

\subsection{Results and progress}
\begin{itemize}
    \item With the optimised energy and energy compressor, the efficiency of configuration 1 increased from $\approx 33\%$ to $\approx 85\%$
    \item The original longitudinal distribution had an energy spread of $6.48\%$ and this has decreased to $5.82\%$
    \item It should also be noted that part of the distribution has been sheared to reduce the tail
    \item Spin tracking is producing very high equilibrium polarization rates, around the asymptotic limit of $92.4\%$ with a polarization time of $\approx 13$ minutes
\end{itemize}

\subsection{Next steps}
\begin{itemize}
    \item Verify spin tracking results and compare with supervisor's code, understand where discrepancies come from
    \item Understand how synchrotron tune plays a role and spin coupling
    \item Explore other cell types, for example TME cells, to try and reduce equilibrium emittance by a factor of 10 
    \item Optimise racetrack configuration
\end{itemize}

\subsection{Skills developed}
\begin{itemize}
    \item Increased understanding of spin motion and tune and how other formalisms introduce radiative depolarizing effects into spin polarization.
\end{itemize}

%---------
% MONTH 5
%---------

\section{Month : dd/mm/2026 - dd/mm/2026}
\subsection{Overview}

\subsection{Work completed}

\begin{itemize}
\item
\end{itemize}

\subsection{Problems encountered}

\textbf{Problem:} text
\begin{itemize}
    \item \textbf{Solution:}
\end{itemize}

\subsection{Results and progress}
\begin{itemize}
    \item 
\end{itemize}

\subsection{Next steps}
\begin{itemize}
    \item 
\end{itemize}

\subsection{Skills developed}
\begin{itemize}
    \item \textbf{Theme:} text
\end{itemize}

\end{document}